% SOURCE_STATEMENT: "Given a Euclidean classical field theory, when does analytic continuation from Euclidean time to real time produce a sensible relativistic quantum field theory?"
% Preparation and final self-audit: CONTRIBUTING.md prompts C and D; source/audit artifacts retained in the local physics_statements_2026_10_11 batch.
% !TeX program = lualatex
\documentclass[11pt]{article}
\usepackage[margin=25mm]{geometry}
\usepackage{fontspec}
\setmainfont{Latin Modern Roman}
\usepackage{amsmath,amssymb,mathtools}
\usepackage{unicode-math}
\setmathfont{Latin Modern Math}
\usepackage{xurl,hyperref}
\hypersetup{colorlinks=true,urlcolor=blue}
\setcounter{secnumdepth}{0}
\setlength{\parindent}{0pt}
\setlength{\parskip}{5pt}
\setlength{\emergencystretch}{3em}
\newcommand{\sref}[2]{\hyperlink{source:#1:#2}{[S#2]}}
\newcommand{\cataloguescope}[1]{\paragraph{Recorded target.}#1}
\newcommand{\R}{\mathbb R}
\newcommand{\C}{\mathbb C}
\newcommand{\N}{\mathbb N}
\newcommand{\Z}{\mathbb Z}
\newcommand{\Q}{\mathbb Q}
\newcommand{\F}{\mathbb F}
\newcommand{\eps}{\varepsilon}
\begin{document}
\section{Euclidean Reconstruction of Relativistic Quantum Field Theories}\label{statement}

\subsection{Definitions and mathematical statement}

\paragraph{Shared notation.}$\mathbb N=\{1,2,\ldots\}$; the standard number fields have their usual meanings. All additional conventions are stated locally.


\paragraph{Definitions and conventions.}
A Euclidean classical field theory is described by Euclidean correlation functions or a probability measure on fields. Reconstruction seeks a positive Hilbert space, vacuum, unitary time evolution and local relativistic observable fields in real time. Reflection positivity is positivity of reflected Euclidean-time correlations and is a crucial compatibility requirement; Euclidean symmetry, regularity and growth conditions also enter reconstruction theorems. Analytic continuation concerns the correlation functions, not merely replacing a real variable by an imaginary one in a formal action.

\paragraph{Target retained from the canonical source catalogue.}
\paragraph{Catalogue wording retained for traceability.}
Given a Euclidean classical field theory, when does analytic continuation from Euclidean time to real time produce a sensible relativistic quantum field theory?

\paragraph{Source scope and qualification.}
The word sensible is a request to identify suitable axioms and analytic conditions. This programme does not say every Euclidean measure admits a Wightman theory or that formal Wick rotation proves reconstruction.
\subsection{Short English statement}
Determine when a Euclidean field model can be reconstructed as a mathematically consistent relativistic quantum theory.
\subsection{Sources}
\begin{itemize}
\item[\hypertarget{source:statement:1}{S1}] ULAM.AI and the UnsolvedMath Contributors, UnsolvedMath, catalogue entry. Catalogue curation under CC BY 4.0; definitions and English exposition added locally.
\url{https://huggingface.co/datasets/ulamai/UnsolvedMath}.
\item[\hypertarget{source:statement:2}{S2}] Reflection Positivity (2017), Oberwolfach Reports, DOI 10.4171/owr/2017/55, pp. 3263--3343. Primary report accessed 2026-10-11; relevant PDF pages 33, 34.
\url{https://doi.org/10.4171/owr/2017/55}.
\item[\hypertarget{source:statement:3}{S3}] Osterwalder--Schrader axioms (catalogue-listed context; not independently consulted in this preparation).
\url{https://doi.org/10.1007/BF01645738}.
\item[\hypertarget{source:statement:4}{S4}] OWR Report 55/2017 (catalogue-listed context; not independently consulted in this preparation).
\url{https://ems.press/content/serial-article-files/46717}.
\end{itemize}
\label{end:statement}
\end{document}
